% =====================================================================
%  references.tex ------ 参考文献（原书 pp.135--140，PDF 155--160）
%  说明：正文的叙述式文献评注照译；文献条目沿用原书英文/日文原貌（不逐条译），
%        OCR 明显讹误处按通行书目改正（如 Stasheff、Twistor 等），并在 README 声明。
% =====================================================================

\chapter*{参考文献}
\addcontentsline{toc}{chapter}{参考文献}
\markboth{参考文献}{参考文献}

% [原文 p.135]
本书所展开的基本想法，多仰赖 Witten 的论文

\noindent [1]\ E. Witten, \textit{Quantum field theory and the Jones polynomial}, Commun. Math. Phys. \textbf{121} (1989), 351--399.

\noindent 此外，作为强调几何视角、概述 Chern--Simons 理论与共形场理论的关系及其向拓扑不变量应用的著作，有

\noindent [2]\ M. F. Atiyah, \textit{The Geometry and Physics of Knots}, Cambridge University Press, 1990.

\noindent 关于量子场论的基础概念，

\noindent [3]\ J. M. Rabin, \textit{Introduction to quantum field theory for mathematicians}, in \textit{Geometry and Quantum Field Theory}, D. S. Freed and K. K. Uhlenbeck, eds., American Mathematical Society, Institute for Advanced Study, 1995, 185--269.

\noindent 有对数学家较为平易的叙述。此外，关于「理论的概要与展望」中所述几何量子化，

\noindent [4]\ N. M. J. Woodhouse, \textit{Geometric Quantization}, Second Edition, Oxford University Press, 1992.

\noindent 有详细论述。关于 Feynman 路径积分等以无限维为对象的几何学的精髓，

\noindent [5]\ 深谷贤治，《無限次元の幾何学》，岩波讲座「现代数学的基础·现代数学的展开 1」，1996.

\noindent 有解说，务请一读。本书所需的向量丛与联络、以及示性类理论，请参照

\noindent [6]\ 森田茂之，《微分形式の幾何学 I, II》，岩波讲座「现代数学的基础」，1996, 1997.

\noindent 等。

% [原文 p.136]
第 1 章以环路群为基础处理 Wess--Zumino--Witten 模型，环路群的中心扩张、环路群上的复线丛等起了重要作用。关于这些，

\noindent [7]\ A. Pressley and G. Segal, \textit{Loop Groups}, Oxford University Press, 1986.

\noindent 有详细叙述。又，作为环路群中心扩张的无穷小版本的概念------仿射 Lie 代数及其表示论，

\noindent [8]\ V. Kac, \textit{Infinite Dimensional Lie Algebras}, Third Edition, Cambridge University Press, 1990.

\noindent 是标准文献。共形场理论由 Belavin--Polyakov--Zamolodchikov [20] 首次作为处理统计力学模型临界现象中无限维对称性的理论登场。进而，从 Wess--Zumino--Witten 模型到 KZ 方程的脉络，在 Knizhnik--Zamolodchikov [33] 中展开。本书中 WZW 模型用环路群上复线丛所作的表述，大体依照 Gawedzki [28]。Riemann 面上的共形场理论，在论文 [45] 中，用仿射 Lie 代数的表示，被表述为模空间上向量丛的理论。

第 2 章叙述共形场理论与拓扑不变量的关系。2.1 节处理的辫群，

\noindent [9]\ J. Birman, \textit{Braids, Links and Mapping Class Groups}, Princeton University Press, 1974.

\noindent 是标准教科书。此书对映射类群也有详细叙述。此外，关于纽结、三维流形中本书未加证明而引用的事实，请参照

\noindent [10]\ 河内明夫 编著，《結び目理論》，Springer-Verlag 东京，1990.

\noindent [11]\ D. Rolfsen, \textit{Knots and Links}, Publish or Perish, Inc., 1976.

\noindent 等。本书处理的 Jones 多项式，在 Jones 的论文 [31] 中从 Iwahori--Hecke 代数表示的观点，连两变量的情形一并定式化。[10] 中也有关于 Jones 多项式的叙述。2.3 节解说的三维流形的 Witten 不变量，据 [43]，与用量子群在 1 的幂根处的表示论所定义者等价。本书替代技术上困难的量子群在 1 的幂根处的表示，尝试了基于共形场理论融合规则的几何展开。

第 3 章涉及的 Vassiliev 理论，

\noindent [12]\ V. A. Vassiliev, \textit{Complements of Discriminants of Smooth Maps: Topology and Applications}, Translations of Mathematical Monographs, American Mathematical Society, 1994.

% [原文 p.137]
\noindent 有详细叙述。3.2 节解说的 Kontsevich 积分依据论文 [36]。作为其原型的 Chen 迭代积分理论，请见 [22]。3.3 节基于 Chern--Simons 微扰理论构造三维流形拓扑不变量，依照 Axelrod--Singer 的论文 [18]。

\noindent 此外，包括正文中已引用者，列出与本书内容关系较深的文献。

\noindent [13]\ V. Arnold, \textit{Topological Invariants of Plane Curves and Caustics}, University Lecture Series 5, American Mathematical Society, 1994.

\noindent [14]\ M. F. Atiyah, \textit{On framings of 3-manifolds}, Topology \textbf{29} (1990), 1--8.

\noindent [15]\ M. F. Atiyah and R. Bott, \textit{The Yang--Mills equation over Riemann surfaces}, Phil. Trans. R. Soc. London A \textbf{308} (1982), 523--615.

\noindent [16]\ M. F. Atiyah, V. K. Patodi and I. M. Singer, \textit{Spectral geometry and Riemannian geometry I}, Math. Proc. Camb. Phil. Soc. \textbf{77} (1975), 43--69; ibid. II ibid. \textbf{78} (1975), 405--432; ibid. III ibid. \textbf{79} (1976), 71--99.

\noindent [17]\ S. Axelrod, S. Della Pietra and E. Witten, \textit{Geometric quantization of Chern--Simons gauge theory}, J. Diff. Geom. \textbf{33} (1991), 782--902.

\noindent [18]\ S. Axelrod and I. M. Singer, \textit{Chern--Simons perturbation theory}, in \textit{Proc. XXth International Conference on Differential Geometric Methods in Theoretical Physics}, S. Catto and A. Rocha eds., World Scientific, 1991, 3--45; ibid. II J. Diff. Geom. \textbf{39} (1994), 173--213.

\noindent [19]\ A. Beauville and Y. Laszlo, \textit{Conformal blocks and generalized theta functions}, Commun. Math. Phys. \textbf{164} (1994), 385--419.

\noindent [20]\ A. A. Belavin, A. M. Polyakov and A. B. Zamolodchikov, \textit{Infinite conformal symmetry in two-dimensional quantum field theory}, Nucl. Phys. \textbf{B241} (1984), 333--380.

\noindent [21]\ R. Bott and C. Taubes, \textit{On the self linking of knots}, J. Math. Phys. \textbf{35} (1994), 5247--5287.

\noindent [22]\ K. T. Chen, \textit{Iterated path integrals}, Bull. Amer. Math. Soc. \textbf{83} (1977), 831--879.

\noindent [23]\ V. G. Drinfel'd, \textit{Quantum groups}, in \textit{Proceedings of the International Congress of Mathematicians, Berkeley}, (1989), 798--820.

% [原文 p.138]
\noindent [24]\ V. G. Drinfel'd, \textit{Quasi-Hopf algebras}, Leningrad Math. J. \textbf{1} (1990), 1419--1457.

\noindent [25]\ D. Evans and Y. Kawahigashi, \textit{Quantum Symmetries on Operator Algebras}, Oxford University Press, 1998.

\noindent [26]\ R. Fenn and C. Rourke, \textit{On Kirby's calculus of links}, Topology \textbf{18} (1979), 1--15.

\noindent [27]\ K. Fukaya, \textit{Morse homotopy and Chern--Simons perturbation theory}, Commun. Math. Phys. \textbf{181} (1996), 37--90.

\noindent [28]\ K. Gawedzki, \textit{Wess--Zumino--Witten conformal field theory}, in \textit{Constructive Quantum Field Theory II}, G. Velo and A. S. Wightman eds., Plenum Press, 1990, 89--120.

\noindent [29]\ E. Guadagnini, M. Martellini and M. Mintchev, \textit{Wilson lines in Chern--Simons theory and link invariants}, Nucl. Phys. \textbf{B330} (1990), 575--607.

\noindent [30]\ M. Jimbo, \textit{A q-difference analogue of $U(\mathfrak{g})$ and the Yang--Baxter equation}, Lett. Math. Phys. \textbf{10} (1985), 63--69.

\noindent [31]\ V. F. R. Jones, \textit{Hecke algebra representations of braid groups and link polynomials}, Ann. Math. \textbf{126} (1987), 335--388.

\noindent [32]\ R. Kirby and P. Melvin, \textit{The 3-manifold invariants of Witten and Reshetikhin--Turaev for $\mathrm{sl}(2,\C)$}, Invent. Math. \textbf{105} (1991), 473--546.

\noindent [33]\ V. G. Knizhnik and A. B. Zamolodchikov, \textit{Current algebra and Wess--Zumino models in two dimensions}, Nucl. Phys. \textbf{B247} (1984), 83--103.

\noindent [34]\ T. Kohno, \textit{Monodromy representations of braid groups and Yang--Baxter equations}, Ann. Inst. Fourier \textbf{37} (1987), 139--160.

\noindent [35]\ T. Kohno, \textit{Topological invariants for 3-manifolds using representations of mapping class groups I}, Topology \textbf{31} (1992), 203--230.

\noindent [36]\ M. Kontsevich, \textit{Vassiliev's knot invariants}, Advances in Soviet Mathematics \textbf{16} (1993), 137--150.

\noindent [37]\ M. Kontsevich, \textit{Intersection theory on the moduli space of curves and the matrix Airy function}, Commun. Math. Phys. \textbf{147} (1992), 1--23.

\noindent [38]\ T. Q. T. Le and J. Murakami, \textit{On Kontsevich's integral for the Homfly polynomial and relations of mixed Euler numbers}, Topology Appl. \textbf{62} (1995), 193--206.

% [原文 p.139]
\noindent [39]\ T. Q. T. Le, J. Murakami and T. Ohtsuki, \textit{On a universal perturbative invariant of 3-manifolds}, Topology \textbf{37} (1998), 539--574.

\noindent [40]\ 村濑元彦，《積分可能系の展開》，岩波讲座「现代数学的基础·现代数学的展开 2」，1997.

\noindent [41]\ H. Murakami, \textit{Quantum $SO(3)$-invariants dominate the $SU(2)$ invariant of Casson and Walker}, Math. Proc. Camb. Phil. Soc. \textbf{117} (1995), 237--249.

\noindent [42]\ T. Ohtsuki, \textit{A polynomial invariant of rational homology 3-spheres}, Invent. Math. \textbf{123} (1996), 241--257.

\noindent [43]\ N. Reshetikhin and V. G. Turaev, \textit{Invariants of 3-manifolds via link polynomials and quantum groups}, Invent. Math. \textbf{103} (1991), 547--597.

\noindent [44]\ A. Tsuchiya and Y. Kanie, \textit{Vertex operators in conformal field theory on $P^{1}$ and monodromy representations of braid groups}, Advanced Studies in Pure Math. \textbf{16} (1988), 297--372.

\noindent [45]\ A. Tsuchiya, K. Ueno and Y. Yamada, \textit{Conformal field theory on universal family of stable curves with gauge symmetries}, Advanced Studies in Pure Math. \textbf{19} (1990), 459--566.

\noindent [46]\ E. Witten, \textit{Two-dimensional gravity and intersection theory on moduli space}, Surveys in Differential Geom. \textbf{1} (1991), 243--310.

\noindent [47]\ T. Yoshida, \textit{An abelianization of $SU(2)$ WZW model}, Ann. Math. \textbf{164} (2006), 1--49.

\bigskip
{\small\noindent 说明：本节的文献评注为原书正文的译文；文献条目按学术惯例保留原书语言（英文/日文），仅对 OCR 明显讹误作最小订正，未逐条核对原书详情。原书 p.141--148 的和文/欧文索引未收入译本，见 README「已知限制」。}
